% ============================================================
\section{Cauchy integrals and the Saffman--Taylor formulation of the Muskat problem}\label{sec:muskat}
% ============================================================

\subsection{Curves and notation}
We identify $\R^2$ with $\C$, $x=x_1+ix_2$. For a vector $w=(w_1,w_2)$ we write $w^\flat:=w_1-iw_2$ (the complex conjugate of $w$ seen as a complex number);
a vector field $w$ is divergence-free and curl-free in an open set if and only if $w^\flat$ is holomorphic there. We write $w^\perp=(-w_2,w_1)$.

\begin{definition}\label{def:curves}
For $k\ge1$ and $\gamma\in(0,1)$ let $\Zc^{k,\gamma}$ be the set of maps $z\in C^{k,\gamma}_{\rm loc}(\R;\R^2)$ such that $z(\alpha)-(\alpha,0)$ is $2\pi$-periodic and
\[
\delta(z):=\inf\Big\{\frac{|\sin\frac{z(\alpha)-z(\beta)}{2}|}{|\sin\frac{\alpha-\beta}{2}|}:\ \alpha,\beta\in\R,\ \alpha-\beta\notin2\pi\mathbb Z\Big\}>0,
\]
where $\sin$ is the complex sine. We norm $\Zc^{k,\gamma}$ by $\|z-(\cdot,0)\|_{C^{k,\gamma}(\T)}$.
\end{definition}

Since $|\sin\frac w2|^2=\sin^2\frac{\Rea w}2+\sinh^2\frac{\Ima w}2$, the condition $\delta(z)>0$ says that $z$ induces an embedding of $\T$ into the cylinder $\T\times\R$,
quantitatively; in particular $|\partial_\alpha z|\ge\delta(z)$. The curve $\Gamma=z(\R)$ is invariant under $x\mapsto x+2\pi e_1$ and $\R^2\setminus\Gamma$ has two connected
components $\Omega^+(z)$ and $\Omega^-(z)$, containing $\{x_2>\max z_2\}$ and $\{x_2<\min z_2\}$ respectively. We write
\[
\tau=\frac{\partial_\alpha z}{|\partial_\alpha z|},\qquad n=\tau^\perp\ \ (\text{so } n=i\tau\text{ as complex numbers}),\qquad \dd s=|\partial_\alpha z|\dd\alpha,
\]
so that $n$ points into $\Omega^+(z)$, and $n_2=\partial_\alpha z_1/|\partial_\alpha z|$. For a function $w$ defined near $\Gamma$ we denote by $w^\pm(\alpha)$ its limits at $z(\alpha)$
from $\Omega^\pm(z)$, by $\jump{w}=w^+-w^-$ its jump and by $\avg{w}=\frac12(w^++w^-)$ its average. The constant $C$ may change from line to line.

\subsection{Periodic Cauchy integrals}

\begin{definition}\label{def:cauchy}
Let $z\in\Zc^{1,\gamma}$ and $f\in C^{0,\gamma}(\T;\R^2)$. The \emph{Cauchy field} of $f$ along $z$ is the vector field $\Cf[z;f]$ on $\R^2\setminus\Gamma$ defined by
\begin{equation}\label{eq:cauchy}
\Cf[z;f]^\flat(x)=\frac1{4\pi i}\int_\T f^\flat(\beta)\,\partial_\beta z(\beta)\,\cot\frac{z(\beta)-x}{2}\dd\beta .
\end{equation}
\end{definition}

Since $\frac12\cot\frac w2=\frac1w+O(w)$, formula \eqref{eq:cauchy} is the periodisation of the classical Cauchy integral
\[
\frac1{2\pi i}\int\frac{f^\flat(\zeta)\dd\zeta}{\zeta-x}.
\]

\begin{lemma}\label{lem:cauchy}
Let $z\in\Zc^{1,\gamma}$, $f\in C^{0,\gamma}(\T;\R^2)$ and $\Cf=\Cf[z;f]$.
\begin{enumerate}[label=\textup{(\alph*)}]
\item $\Cf$ is real analytic and $2\pi$-periodic in $x_1$ on $\R^2\setminus\Gamma$, with $\dvg\Cf=\curl\Cf=0$ there.
\item \textup{(Plemelj--Privalov)} The restrictions of $\Cf$ to $\Omega^\pm(z)$ extend to $C^{0,\gamma}(\overline{\Omega^\pm(z)})$ with
\[
\|\Cf\|_{C^{0,\gamma}(\overline{\Omega^\pm})}\le C\|f\|_{C^{0,\gamma}},
\]
where $C$ depends on $\gamma$, $\delta(z)$ and $\|z\|_{\Zc^{1,\gamma}}$. Moreover
$\jump{\Cf}=f$, and $\avg{\Cf}$ is given by \eqref{eq:cauchy} with $x=z(\alpha)$ and the integral understood in the principal value sense.
\item As $x_2\to\pm\infty$, $\Cf(x)\to\Cf_{\pm\infty}$ uniformly in $x_1$, where $\Cf_{\pm\infty}^\flat=\pm\frac1{4\pi}\int_\T f^\flat\partial_\beta z\dd\beta$.
In particular, for a constant vector $c$ we have $\Cf[z;c]=\pm\frac c2$ in $\Omega^\pm(z)$, and for a scalar $\phi$
\begin{equation}\label{eq:limit-source}
\Cf[z;\phi n]_{\pm\infty}=\pm\frac1{4\pi}\Big(\int_\T\phi\,|\partial_\beta z|\dd\beta\Big)\ee .
\end{equation}
\item If $j\ge0$, $z\in\Zc^{j+1,\gamma}$ and $f\in C^{j,\gamma}$, then $\Cf\in C^{j,\gamma}(\overline{\Omega^\pm(z)})$ and $\Cf^\pm\in C^{j,\gamma}(\T)$, with bounds.
\item \textup{(Representation)} Let $w$ be a bounded vector field on $\R^2\setminus\Gamma$, $2\pi$-periodic in $x_1$, divergence-free and curl-free in $\R^2\setminus\Gamma$, whose restrictions
to $\Omega^\pm(z)$ extend continuously to the closures with $\jump{w}\in C^{0,\gamma}(\T)$, and with $w(x)\to0$ as $x_2\to\pm\infty$. Then $w=\Cf[z;\jump{w}]$.
\end{enumerate}
\end{lemma}

\begin{proof}
(a) is immediate. (b) and (d) are the periodic versions of the Plemelj--Privalov theorem and are discussed in Appendix~\ref{app:cauchy}.
(c) When $x_2\to+\infty$ (resp.\ $-\infty$), $\cot\frac{z(\beta)-x}2\to i$ (resp.\ $-i$) uniformly in $\beta$, which gives $\Cf_{\pm\infty}$. If $f=c$ is constant, then
$\partial_x\int_\T c^\flat\partial_\beta z\cot\frac{z(\beta)-x}2\dd\beta=-c^\flat\int_\T\partial_\beta\big[\cot\frac{z(\beta)-x}{2}\big]\dd\beta=0$ by periodicity, so $\Cf[z;c]$ is constant
in each component and equals its limit $\pm\frac1{4\pi}c^\flat\cdot2\pi=\pm\frac{c^\flat}{2}$; recall $\int_\T\partial_\beta z\dd\beta=2\pi$.
For $f=\phi n$ we have $n^\flat\partial_\beta z=-i|\partial_\beta z|$, so $\Cf_{+\infty}^\flat=-\frac i{4\pi}\int\phi|\partial_\beta z|$, i.e.\ \eqref{eq:limit-source}.
(e) The field $W=w-\Cf[z;\jump w]$ is bounded, periodic, continuous across $\Gamma$, and $W^\flat$ is holomorphic off $\Gamma$; hence $W^\flat$ is entire (Morera), and by Liouville's
theorem on the cylinder it is constant. By (c), $W\to-\Cf_{+\infty}$ as $x_2\to+\infty$ and $W\to-\Cf_{-\infty}=\Cf_{+\infty}$ as $x_2\to-\infty$. Therefore $\Cf_{+\infty}=0$ and $W\equiv0$.
\end{proof}

\subsection{The adjoint Neumann--Poincar\'e operator}

\begin{definition}\label{def:K}
For $z\in\Zc^{2,\gamma}$ and a real function $\phi\in C^{0,\gamma}(\T)$ we set
\[
\Kop(z)\phi:=\avg{\Cf[z;\phi n]}\cdot n .
\]
\end{definition}
Since the jump of $\Cf[z;\phi n]$ is normal, Lemma~\ref{lem:cauchy}(b) gives
\begin{equation}\label{eq:K-traces}
\Cf[z;\phi n]^\pm\cdot n=\Kop(z)\phi\pm\tfrac12\phi,\qquad \Cf[z;\phi n]^+\cdot\tau=\Cf[z;\phi n]^-\cdot\tau .
\end{equation}
In potential-theoretic terms, $\Cf[z;\phi n]$ is the gradient of the single layer potential of $\phi$ and $\Kop(z)$ is the (periodic) adjoint Neumann--Poincar\'e operator.

Lemma~\ref{lem:K}(a) and Proposition~\ref{prop:invert} below are H\"older-space versions of classical results of C\'ordoba, C\'ordoba and Gancedo \cite{CCG11} in the horizontally
periodic setting: the smoothing property of $\Kop(z)$ (cf.\ \cite[Lemmas~3.1 and~5.2]{CCG11}) and, in Sobolev spaces, the invertibility of $I-2\xi\Kop(z)$ for $|\xi|\le1$, by a
harmonic-energy argument (cf.\ \cite[Sec.~4]{CCG11}). They are recalled in \cite[Secs.~3.1.3, 8.2 and~8.3]{LM26}, whose operator $P$ corresponds to our $2\Kop$, and the versions
for asymptotically flat curves are \cite[Props.~6.1 and~8.3]{LM26}.

\begin{lemma}\label{lem:K}
Let $k\ge2$ and $z\in\Zc^{k,\gamma}$.
\begin{enumerate}[label=\textup{(\alph*)}]
\item $\Kop(z)\phi(\alpha)=\int_\T k_z(\alpha,\beta)\phi(\beta)\dd\beta$, with
\[
k_z(\alpha,\beta)=\frac{|\partial_\beta z(\beta)|}{4\pi|\partial_\alpha z(\alpha)|}\,\Ima\Big[\partial_\alpha z(\alpha)\cot\frac{z(\beta)-z(\alpha)}{2}\Big],
\]
and $k_z$ extends from $\{\alpha\neq\beta\}$ to a function in $C^{k-2,\gamma}(\T\times\T)$. Consequently $\Kop(z)$ maps $C^0(\T)$ boundedly into $C^{k-2,\gamma}(\T)$, and for $k\ge3$ it is
compact on $C^{j,\gamma}(\T)$ for $0\le j\le k-3$.
\item $\int_\T\big(\Kop(z)\phi\big)\,|\partial_\alpha z|\dd\alpha=0$ for every $\phi\in C^{0,\gamma}(\T)$.
\end{enumerate}
\end{lemma}

\begin{proof}
(a) The complex density of $\phi n$ is $\phi\, n^\flat\partial_\beta z=-i\phi|\partial_\beta z|$. Using $w\cdot n=\Rea(w^\flat n)$ and $n=i\partial_\alpha z/|\partial_\alpha z|$, the principal value formula of
Lemma~\ref{lem:cauchy}(b) gives the stated kernel, once we show it is bounded (then no principal value is needed). Away from the diagonal $\{\alpha-\beta\in2\pi\mathbb Z\}$ the kernel is
$C^{k-1,\gamma}$, because $\delta(z)>0$. Near the diagonal write $\cot\frac w2=\frac2w+g(w)$ with $g$ analytic for $|w|<2\pi$, and
\[
z(\beta)-z(\alpha)=(\beta-\alpha)Z_1(\alpha,\beta),\qquad \partial_\alpha z(\alpha)-Z_1(\alpha,\beta)=-(\beta-\alpha)Z_2(\alpha,\beta),
\]
with $Z_1=\int_0^1\partial_\alpha z(\alpha+s(\beta-\alpha))\dd s\in C^{k-1,\gamma}$, $Z_2=\int_0^1\!\!\int_0^1 s\,\partial^2_\alpha z(\alpha+rs(\beta-\alpha))\dd r\dd s\in C^{k-2,\gamma}$ and $|Z_1|\ge c>0$. Then
\[
\Ima\Big[\frac{2\partial_\alpha z(\alpha)}{z(\beta)-z(\alpha)}\Big]=\frac{2}{\beta-\alpha}\Ima\Big[\frac{\partial_\alpha z(\alpha)-Z_1}{Z_1}\Big]=-2\,\Ima\Big[\frac{Z_2}{Z_1}\Big]\in C^{k-2,\gamma},
\]
while $\partial_\alpha z(\alpha)g(z(\beta)-z(\alpha))\in C^{k-1,\gamma}$. Differentiating under the integral sign gives $\Kop(z):C^0\to C^{k-2,\gamma}$, and compactness follows from the compact
embedding $C^{k-2,\gamma}(\T)\subset C^{j,\gamma}(\T)$ for $j\le k-3$.

(b) Let $w=\Cf[z;\phi n]$. Since $w$ is divergence-free and periodic, its flux through $\Gamma$ from above equals its flux through $\{x_2=L\}$ for every large $L$, and letting $L\to\infty$ we get
from \eqref{eq:limit-source} $\int_\T w^+\cdot n\,\dd s=2\pi(\Cf_{+\infty})_2=\frac12\int\phi\dd s$. Similarly $\int_\T w^-\cdot n\dd s=2\pi(\Cf_{-\infty})_2=-\frac12\int\phi\dd s$. Adding,
$\int\avg{w}\cdot n\dd s=0$.
\end{proof}

\begin{proposition}\label{prop:invert}
Let $|A|<1$, $k\ge3$ and $z\in\Zc^{k,\gamma}$. Then $I-2A\Kop(z)$ is an isomorphism of $C^{j,\gamma}(\T)$ for $0\le j\le k-3$. Moreover, if $f\in C^{j,\gamma}(\T)$ with $0\le j\le k-2$, the
solution $\sig\in C^{0,\gamma}(\T)$ of $(I-2A\Kop(z))\sig=f$ belongs to $C^{j,\gamma}(\T)$, with $\|\sig\|_{C^{j,\gamma}}\le C\|f\|_{C^{j,\gamma}}$, and $\int\sig\dd s=\int f\dd s$.
\end{proposition}

\begin{proof}
By Lemma~\ref{lem:K}(a) and the Riesz--Schauder theory, $I-2A\Kop(z)$ is a Fredholm operator of index zero on $C^{j,\gamma}(\T)$, $0\le j\le k-3$; so it suffices to prove injectivity on $C^{0,\gamma}$.
Let $\phi\in C^{0,\gamma}$ with $\phi=2A\Kop(z)\phi$. If $A=0$ there is nothing to prove, so let $A\ne0$. By Lemma~\ref{lem:K}(b), $\int\phi\dd s=0$, so $w:=\Cf[z;\phi n]$ tends to $0$ as
$x_2\to\pm\infty$, by \eqref{eq:limit-source}; being the conjugate of a periodic holomorphic function, it does so exponentially fast. Since the jump of $w$ is normal, $w$ is curl-free in $\mathcal D'(\R^2)$,
and since its circulation along horizontal lines tends to zero, $w=\nabla\Phi$ for a periodic, Lipschitz function $\Phi$ which is $C^{1,\gamma}$ up to $\Gamma$ from both sides and bounded.
By \eqref{eq:K-traces} and $\Kop(z)\phi=\phi/(2A)$,
\[
\frac{w^+\cdot n}{1+A}=\frac{w^-\cdot n}{1-A}=\frac{\phi}{2A}=:g .
\]
Integrating by parts on one period of $\Omega^\pm\cap\{|x_2|<L\}$ and letting $L\to\infty$ (the contributions of the horizontal boundaries vanish by the exponential decay, the lateral ones cancel by periodicity),
\[
\frac1{1+A}\int_{\Omega^+}|w|^2+\frac1{1-A}\int_{\Omega^-}|w|^2=-\int_\T\Phi\,\frac{w^+\cdot n}{1+A}\dd s+\int_\T\Phi\,\frac{w^-\cdot n}{1-A}\dd s=0,
\]
where the integrals over $\Omega^\pm$ are over one period. Hence $w\equiv0$ and $\phi=\jump{w}\cdot n=0$.

For the regularity statement with $j=k-2$: the solution $\sig\in C^{k-3,\gamma}$ satisfies $\sig=f+2A\Kop(z)\sig$, and $\Kop(z)\sig\in C^{k-2,\gamma}$ by Lemma~\ref{lem:K}(a). Finally
$\int\sig\dd s=\int f\dd s$ by Lemma~\ref{lem:K}(b).
\end{proof}

\subsection{The Muskat problem with viscosity jump}
Let $z\in\Zc^{k,\gamma}$, $k\ge3$, and $\rho=\pm1$ in $\Omega^\pm(z)$. The \emph{instantaneous Darcy problem} consists in finding $u$ and $p$, $C^{0,\gamma}$ and $C^{1,\gamma}$ up to $\Gamma$ from
both sides, such that
\begin{equation}\label{eq:darcy-problem}
\dvg u=0,\quad (1\pm A)\,u=-\nabla p\mp\ee\ \text{ in }\Omega^\pm(z),\qquad \jump{p}=0,\ \ \jump{u\cdot n}=0\ \text{ on }\Gamma,
\end{equation}
with $u\to0$ and $\nabla p+\rho\ee\to0$ as $x_2\to\pm\infty$ (the fluids are at rest at infinity). We introduce the \emph{Darcy field} and the \emph{Saffman--Taylor function}
\[
v:=\mu(\rho)u=-\nabla p-\rho\ee,\qquad \sig:=\jump{\partial_n p}=\partial_np^+-\partial_np^-\ \text{ on }\Gamma .
\]
Proposition~\ref{prop:sigma} and Corollary~\ref{cor:muskat} below are the Saffman--Taylor formulation $(M\sigma)$ of Luna-Velasco and Mengual \cite[Sec.~3.1]{LM26}, transcribed to the
horizontally periodic setting: \eqref{eq:v-sigma} says that $\nabla p=\Cf[z;\sig n]$, which is the representation \cite[(3.2)]{LM26} of the pressure gradient; \eqref{eq:U-sigma} is
\cite[(3.6)]{LM26}; the identity $\sig=-2(n_2+AU)$ is \cite[(3.7a) and (3.8)]{LM26}; and \eqref{eq:sigma-eq} is \cite[(5.12)]{LM26}. The dictionary is $\jump\theta=-2$ for the phase jump
of \cite{LM26}, whose Saffman--Taylor density is $|\partial_\alpha z|\sig$ and whose operator $P$ corresponds to $2\Kop(z)$. The proof below follows the derivation of $(M\sigma)$ in
\cite[Sec.~3.1.5]{LM26}.

\begin{proposition}\label{prop:sigma}
Let $|A|<1$, $k\ge3$, $z\in\Zc^{k,\gamma}$, and let $\sig\in C^{k-2,\gamma}(\T)$ be the solution of
\begin{equation}\label{eq:sigma-eq}
\big(I-2A\Kop(z)\big)\sig=-2n_2 .
\end{equation}
Then \eqref{eq:darcy-problem} has a solution, unique up to an additive constant in $p$, given by
\begin{equation}\label{eq:v-sigma}
v=-\Cf[z;\sig n]\mp\ee\quad\text{and}\quad u=\frac{v}{1\pm A}\qquad\text{in }\Omega^\pm(z).
\end{equation}
Its Saffman--Taylor function is $\sig$, the normal velocity $U:=u\cdot n$ on $\Gamma$ is
\begin{equation}\label{eq:U-sigma}
U=-\Kop(z)\sig,
\end{equation}
and moreover $\sig=-2(n_2+AU)$ and $\int_\T\sig\dd s=-4\pi$.
\end{proposition}

\begin{proof}
Define $v$ and $u$ by \eqref{eq:v-sigma}. By \eqref{eq:K-traces}, $v^\pm\cdot n=-\Kop\sig\mp\frac12\sig\mp n_2$, and \eqref{eq:sigma-eq} says $\frac12\sig+n_2=A\Kop\sig$. Therefore
$v^\pm\cdot n=-(1\pm A)\Kop\sig$, i.e.\ $u^+\cdot n=u^-\cdot n=-\Kop\sig$, which gives $\jump{u\cdot n}=0$ and \eqref{eq:U-sigma}; then $\sig=-2n_2+2A\Kop\sig=-2(n_2+AU)$.
Integrating \eqref{eq:sigma-eq} and using Lemma~\ref{lem:K}(b) and $\int n_2\dd s=\int\partial_\alpha z_1\dd\alpha=2\pi$, we get $\int\sig\dd s=-4\pi$; by \eqref{eq:limit-source},
$v\to\mp\frac1{4\pi}(\int\sig\dd s)\ee\mp\ee=0$ as $x_2\to\pm\infty$. Next, $v+\rho\ee=-\Cf[z;\sig n]$ in $\R^2\setminus\Gamma$; this field is curl-free in $\mathcal D'(\R^2)$, because its jump is normal,
and its circulation along horizontal lines vanishes (its limits at $\pm\infty$ are vertical). Hence $v+\rho\ee=-\nabla p$ for a continuous function $p$, periodic in $x_1$, and
$\partial_np^\pm=\Cf[z;\sig n]^\pm\cdot n=\Kop\sig\pm\frac12\sig$, so that $\jump{\partial_np}=\sig$. The field $u$ is divergence-free in each $\Omega^\pm$ and has no normal jump, so $\dvg u=0$ in
$\mathcal D'(\R^2)$. Regularity follows from Lemma~\ref{lem:cauchy}(d).

For uniqueness, let $(u',p')$ be the difference of two solutions. Then $\dvg u'=0$, $(1\pm A)u'=-\nabla p'$ in $\Omega^\pm$, $p'$ is continuous, $u'\cdot n$ has no jump, and $u'\to0$ at infinity,
exponentially fast since $u'$ is the gradient of a periodic harmonic function in $\{|x_2|>L\}$. Integrating by parts over one period,
\[
\int(1+A\rho)|u'|^2=-\int\nabla p'\cdot u'=\int_\T p'\,\big(u'^+-u'^-\big)\cdot n\dd s=0,
\]
so $u'=0$ and $\nabla p'=0$.
\end{proof}

The identity $\int_\T\sig\dd s=-4\pi$ is the periodic counterpart of the normalisation $\sig^{(0)}+2\in H^{k-1}$ of \cite[Prop.~6.3]{LM26}. It also follows by integrating
$\sig=-2(n_2+AU)$ over a period, since $\int_\T U\dd s=0$ by Lemma~\ref{lem:K}(b) (conservation of area).

Since the interface moves with the normal velocity of the fluid, Proposition~\ref{prop:sigma} gives the following formulation of the Muskat problem with viscosity jump, in which the only
operator to be inverted acts on the Saffman--Taylor function.

\begin{corollary}\label{cor:muskat}
A family of curves $z(\cdot,t)\in\Zc^{k,\gamma}$, $k\ge3$, solves the Muskat problem with viscosity jump if and only if
\[
\partial_t z\cdot n=-\Kop(z)\sig,\qquad \big(I-2A\Kop(z)\big)\sig=-2n_2 .
\]
\end{corollary}

\begin{proof}
By definition, $z$ solves the Muskat problem if, for every $t$, $\partial_tz\cdot n=u\cdot n$ on $\Gamma(t)$ (the kinematic condition, cf.\ \cite[(3.1)]{LM26}), where $u$ is the velocity
of the solution of \eqref{eq:darcy-problem} for the curve $z(\cdot,t)$. By Proposition~\ref{prop:sigma}, $u\cdot n=-\Kop(z)\sig$, where $\sig$ is the unique solution of
\eqref{eq:sigma-eq} (Proposition~\ref{prop:invert}).
\end{proof}

\begin{remark}[Duality with the vorticity formulation]\label{rem:duality}
For a real function $\gamma$, the complex densities of $\gamma\tau$ and $\gamma n$ are $\gamma|\partial_\beta z|$ and $-i\gamma|\partial_\beta z|$, so $\Cf[z;\gamma\tau]^\flat=i\,\Cf[z;\gamma n]^\flat$, i.e.\
$\Cf[z;\gamma\tau]$ is obtained by rotating $\Cf[z;\gamma n]$ by $-\pi/2$; this is the relation between the representations of the velocity and of the pressure gradient in \cite[(3.2)]{LM26}.
In particular $\avg{\Cf[z;\gamma\tau]}\cdot\tau=\avg{\Cf[z;\gamma n]}\cdot n=\Kop(z)\gamma$. The velocity is curl-free off $\Gamma$ with jump $\gamma\tau$, $\gamma=\jump{u\cdot\tau}$, hence
$u=\Cf[z;\gamma\tau]$ by Lemma~\ref{lem:cauchy}(e). Hence $\avg u=\mathrm{BR}(z,\varpi)$, where
$\mathrm{BR}(z,\varpi)^\flat(\alpha):=\frac1{4\pi i}\pv\int_\T\varpi(\beta)\cot\frac{z(\alpha)-z(\beta)}2\dd\beta$ is the (periodised) Birkhoff--Rott integral and $\varpi:=-\gamma|\partial_\alpha z|$
is the vorticity density of the sheet per unit $\alpha$ (so that $\curl u=-\gamma\,\mathcal H^1\llcorner\Gamma$; this is the density $\varpi$ of \cite{LM26}). The tangential part of Darcy's law,
$\jump{v\cdot\tau}=-\jump{\rho}\tau_2=-2\tau_2$, reads $\gamma+2A\avg u\cdot\tau=-2\tau_2$, i.e.\ $(I+2A\Kop(z))\gamma=-2\tau_2$; in terms of $\varpi$ this is
$\varpi-2A\,\mathrm{BR}(z,\varpi)\cdot\partial_\alpha z=2\partial_\alpha z_2$, which is $(I+AP)\varpi=-\jump\theta\partial_\alpha z_2$ of \cite[Sec.~3.1.3]{LM26} with $\jump\theta=-2$, and the
equation for the vorticity of \cite{CCG11} up to the sign conventions used there for the jumps and the Atwood number (cf.\ \cite[Sec.~3.1.3]{LM26}). Thus the vorticity is obtained by
inverting $I+2A\Kop(z)$ and the Saffman--Taylor function by inverting $I-2A\Kop(z)$: as observed in \cite[Sec.~3.1.3]{LM26}, both problems reduce to inverting operators of the form $I\pm AP$,
and both are covered by Proposition~\ref{prop:invert}. In the vorticity formulation the normal velocity $U=\avg u\cdot n=-\avg{\Cf[z;\gamma n]}\cdot\tau$ is a singular integral of $\gamma$,
whereas in \eqref{eq:U-sigma} it is the smoothing operator $\Kop(z)$ applied to $\sig$. This concerns only the map from the density to the normal velocity at a fixed interface (cf.\
\cite[Prop.~8.3]{LM26}) and is not a gain of regularity for the evolution: the loss of derivatives behind the Saffman--Taylor instability sits in the dependence of $\Kop(z)$ on $z$, and the
singular operator $\Top(z)$ of Section~\ref{x:vortex-sheet} reappears when $\gamma$ is recovered from $\sig$ (Proposition~\ref{prop:gamma}). As explained in \cite[Secs.~3.1.4 and~3.3.5]{LM26},
the advantage of the Saffman--Taylor formulation appears for subsolutions with a mixing zone: the pressure remains harmonic in each region, whereas the velocity carries vorticity in the
interior of the mixing zone, so that a vorticity formulation would couple the reconstruction of the velocity to the interior extension of the stream function of the relaxed flux ($\Psi$ in
Section~\ref{sec:ansatz}, the potential $g$ in \cite{LM26}); in the strips of Section~\ref{sec:ansatz}, $\curl\bar u=-A\curl\bar m=-\frac{A(1-\rhom^2)}{2(1-A^2)}\Delta\Psi$ in $S_t$.
\end{remark}

\subsection{The vortex-sheet strength and the pinch set}\label{x:vortex-sheet}
The letter $\gamma$ denotes both the H\"older exponent and the vortex-sheet strength; the meaning is always clear from the context.
For $z\in\Zc^{2,\gamma}$ and a real function $\phi\in C^{0,\gamma}(\T)$ we also set
\[
\Top(z)\phi:=\avg{\Cf[z;\phi n]}\cdot\tau .
\]
Unlike $\Kop(z)$, the operator $\Top(z)$ is a singular integral operator: the argument in the proof of Lemma~\ref{lem:K}(a) shows that, for $z\in\Zc^{k,\gamma}$ with $k\ge2$,
$\Top(z)\phi-\frac{1}{2|\partial_\alpha z|}H\big[|\partial_\alpha z|\phi\big]$ is given by a kernel in $C^{k-2,\gamma}(\T\times\T)$, where
$H\phi(\alpha):=\frac1{2\pi}\pv\int_\T\phi(\beta)\cot\frac{\alpha-\beta}2\dd\beta$ is the periodic Hilbert transform; for the flat interface $\Top=\frac12H$ (cf.\ \cite[(8.14)]{LM26} on $\R$).
The operator $\Top(z)$ is bounded on $C^{j,\gamma}(\T)$ for $z\in\Zc^{j+1,\gamma}$ by Lemma~\ref{lem:cauchy}(d). As in \cite[Sec.~3.1.3]{LM26}, the vortex-sheet strength of the Muskat
problem can be recovered from $\sig$ by a local formula in $\sig$ and $\Top(z)\sig$:

\begin{proposition}\label{prop:gamma}
In the setting of Proposition~\ref{prop:sigma}, the vortex-sheet strength $\gamma=\jump{u\cdot\tau}$ is
\begin{equation}\label{eq:gamma-sigma}
\gamma=-\frac{2\big(\tau_2-A\,\Top(z)\sig\big)}{1-A^2},
\end{equation}
and it is the unique solution in $C^{0,\gamma}(\T)$ of $(I+2A\Kop(z))\gamma=-2\tau_2$. In particular $\gamma\in C^{k-2,\gamma}(\T)$.
\end{proposition}

\begin{proof}
By \eqref{eq:v-sigma}, $v^\pm\cdot\tau=-\Top(z)\sig\mp\tau_2$ (the tangential traces of $\Cf[z;\sig n]$ coincide). Since $u^\pm=v^\pm/(1\pm A)$,
\[
\gamma=\frac{-\Top\sig-\tau_2}{1+A}-\frac{-\Top\sig+\tau_2}{1-A}=\frac{2A\,\Top\sig-2\tau_2}{1-A^2}.
\]
The velocity is curl-free and divergence-free off $\Gamma$, has no normal jump, tends to $0$ at infinity and has tangential jump $\gamma$, so $u=\Cf[z;\gamma\tau]$ by Lemma~\ref{lem:cauchy}(e);
Remark~\ref{rem:duality} then gives $(I+2A\Kop(z))\gamma=-2\tau_2$, and uniqueness and regularity follow from Proposition~\ref{prop:invert} applied with $-A$ in place of $A$.
\end{proof}

Identity \eqref{eq:gamma-sigma} is the left relation of \cite[(3.7b)]{LM26}, $(1-A^2)\varpi=\jump\theta(AQ-\ee)\cdot\partial_\alpha z$, in our normalisation: $\jump\theta=-2$,
$\varpi=-|\partial_\alpha z|\gamma$ (Remark~\ref{rem:duality}), the Saffman--Taylor density of \cite{LM26} is $|\partial_\alpha z|\sig$, and their mean pressure-gradient trace is
$Q=\avg{\Cf[z;\sig n]}$, so that $Q\cdot\partial_\alpha z=|\partial_\alpha z|\,\Top(z)\sig$ and $\ee\cdot\partial_\alpha z=|\partial_\alpha z|\tau_2$. Luna-Velasco and Mengual \cite[Sec.~3.1.3]{LM26}
recover $\varpi$ from $\sig$ in exactly this way, and the equation $(I+2A\Kop(z))\gamma=-2\tau_2$ is their $(I+AP)\varpi=-\jump\theta\partial_\alpha z_2$; only the uniqueness and regularity
statements in periodic H\"older spaces are added here.

\begin{definition}\label{def:unstable}
Let $z_0\in\Zc^{k,\gamma}$, $k\ge3$, and $\Kop_0:=\Kop(z_0)$, $\Top_0:=\Top(z_0)$. The Saffman--Taylor function, the normal velocity and the vortex-sheet strength of the Muskat data
$\rho_0=\mathbf 1_{\Omega^+(z_0)}-\mathbf 1_{\Omega^-(z_0)}$ are
\[
\sig_0:=-2\big(I-2A\Kop_0\big)^{-1}[n_{0,2}],\qquad U_0:=-\Kop_0\sig_0,\qquad \gamma_0:=-2\big(I+2A\Kop_0\big)^{-1}[\tau_{0,2}] .
\]
The data are \emph{fully unstable} if $\sig_0<0$ on $\T$. The \emph{pinch set} is $\Pinch(z_0):=\{\alpha:\ \sig_0(\alpha)\ge0,\ \gamma_0(\alpha)=0\}$.
\end{definition}

The functions $\sig_0$ and $U_0$ are the periodic counterparts of $\sig^{(0)}$ and of the normal component of $z^{(1)}$ in \cite[(5.12)--(5.13)]{LM26}, and full instability is the
hypothesis of \cite[Thm.~1.1]{LM26}. For the relation of $\Pinch(z_0)$ with the pinch singularity of \cite{Men22} see Remark~\ref{rem:hull}(b).

\begin{example}\label{ex:A0}
For the flat interface $z_0(\alpha)=(\alpha,0)$ the kernel $k_{z_0}$ vanishes (the cotangent of a real number is real), so $\Kop_0=0$, $\sig_0\equiv-2$, $U_0\equiv0$ and $\gamma_0\equiv0$.
If $A=0$, then $\sig_0=-2n_{0,2}$ and $\gamma_0=-2\tau_{0,2}$ for every $z_0$: the data are fully unstable if and only if $z_0$ is a graph, and $\Pinch(z_0)=\{n_0=-\ee\}=\{\tau_0=-e_1\}$.
These formulas, and the resulting characterisation of fully unstable data for $A=0$, are those of \cite[Secs.~1.2 and~3.1.3]{LM26}: there $\sig=\jump\theta$ for the flat interface, and
$\sig=\jump\theta\,\partial_\alpha z_1$, $\varpi=-\jump\theta\,\partial_\alpha z_2$ for $A=0$, with $\jump\theta=-2$.
\end{example}

For $z_0\in\Zc^{6,\gamma}$, Propositions~\ref{prop:invert} (with $k=6$, $j=4$) and \ref{prop:gamma}, and Lemma~\ref{lem:K}(a), give $\sig_0,U_0,\gamma_0\in C^{4,\gamma}(\T)$. The quantities
$\Phi$ and $\kmax$ of \eqref{eq:Phi} are then well defined, $\kmax=0$ exactly on $\Pinch(z_0)$, and since the expression under the square root can only vanish where
$\sig_0=\gamma_0=0$, we have $\kmax\in C^{4,\gamma}(\T)$ whenever $\Pinch(z_0)=\emptyset$.
